\documentclass[11pt,letterpaper]{article}
\usepackage[margin=1in]{geometry}
\usepackage[T1]{fontenc}
\usepackage{lmodern,microtype,amsmath,amssymb,amsthm,mathtools}
\usepackage[hidelinks]{hyperref}
\setlength{\emergencystretch}{2em}
\numberwithin{equation}{section}
\newtheorem{theorem}{Theorem}[section]
\newtheorem{proposition}[theorem]{Proposition}
\newtheorem{lemma}[theorem]{Lemma}
\newtheorem{corollary}[theorem]{Corollary}
\DeclareMathOperator{\lk}{lk}
\newcommand{\HB}[1]{\mathcal H_{#1}}
\newcommand{\TB}[1]{\mathcal T_{#1}}
\title{Recognition of smooth homotopy four-spheres\\from marked Morse heights}
\author{Edward B. Baker III}
\date{Working manuscript, 25 November 2026}

\begin{document}
\maketitle
\begin{abstract}
Every Morse function on a nonstandard smooth homotopy four-sphere has
at least two index-two points. For an immersion into Euclidean
five-space, the indexed normal area and an integrated Weyl estimate
give smooth recognition at chiral Weyl energy at most \(64\pi^2/9\).
Below \(8\pi^2\), a putative nonstandard source has exactly two
index-two points in more than three quarters of height directions.
We carry the framed attaching tubes of such a height through the
two end-handlebody reductions. In the one-sided case an unknotted
component or tunnel number one in the actual two-component surgery
link suffices. In the mixed case either of two intermediate surgery
manifolds, \(S^3\) or \(\#_2(S^1\times S^2)\), suffices. The proofs
retain the boundary identifications. Selection of a diagram satisfying
one of these conditions remains a separate geometric problem.
\end{abstract}

\section{Morse functions with few middle handles}\label{sec:marked}

All manifolds and maps are smooth. Let \(M\) be a closed connected
oriented homotopy four-sphere. A diffeomorphism to the standard sphere
can be made orientation preserving by reflection of the target.
For a Morse function \(h\) on \(M\)
write \(n_i\) for its index-\(i\) count. Here Morse means that all
critical points are nondegenerate, with coincident critical values
allowed. An index-\(i\) handle is
\(D^i\times D^{4-i}\), attached by its framed tube
\(S^{i-1}\times D^{4-i}\). Its belt sphere is
\(\{0\}\times S^{3-i}\) in the outgoing boundary. Reversal
interchanges the attaching and belt data and changes index \(i\) to
\(4-i\). Write
\[
 \HB{k}=\natural_k(S^1\times D^3),\qquad
 \TB{k}=\natural_k(S^2\times D^2),\qquad
 Y_k=\#_k(S^1\times S^2),
\]
where \(\natural\) denotes boundary connected sum. In particular,
\(\HB{0}=D^4\) and \(Y_0=S^3\).

\begin{lemma}[Marked end reductions]\label{lem:trees}
The handles of \(h\) yield an ordered presentation of the actual \(M\)
with reduced index counts
\[
 (1,\ell,n_2,r,1),\qquad
 \ell=n_1-n_0+1,\quad r=n_3-n_4+1,\quad n_2=\ell+r .
\]
The lower and upper dual \(0\)--\(1\) bodies are \(\HB{\ell}\) and
\(\HB{r}\). Their identifications with the original end bodies
transport the middle attaching tubes and the outgoing gluing.
\end{lemma}

\begin{proof}
To separate tied critical values when needed, use small disjoint
scalar bumps constant near the individual critical points. The
supports of the derivatives avoid the critical points, so on each
such compact support the differential of \(h\) is bounded away
from zero. Small amplitudes keep the same critical points and
Hessians. This makes an auxiliary scalar function on \(M\).

Use the final rearrangement theorem and its self-indexing formulation
to order the handles by index on \(M\)
\cite[Theorems 4.7--4.8]{Milnor1965HCobordism}. These results keep the critical
points and their indices, and one can choose disjoint attaching tubes
at a common incoming level of the index-two block. The ordered
decomposition uses a gradient-like field and regular collars. It is
an auxiliary Morse presentation if the original critical values do
not have this order. The critical count belongs to \(h\). Separating
ties, changing the order within or between index blocks, and sliding
framed handles can produce cuts of auxiliary presentations. To use
the geometry of an unmodified level of \(h\), one needs a comparison
with the particular cut being used. Within a block of disjoint
attachments one may attach the handles in any order on the common
level, using a product collar to identify the full outgoing boundary.

Take collars on the cuts immediately before and after the index-two
block. The lower body \(P\) is connected because subsequent
attaching regions have connected attaching sphere and so cannot join
distinct components. The same argument applied to the reversed
decomposition makes the upper dual body \(Q\) connected. Form a graph
in each body with a vertex for every \(0\)-handle and an edge for
every \(1\)-handle. Choose a spanning tree in each graph. Attach
the edges of the tree first. At any leaf, the leaf ball and the
edge to the rest amount to an enlargement across a boundary disk.
After rounding corners the enlargement is a ball. Equivalently the
attaching zero-sphere of the edge meets the belt of the leaf
\(0\)-handle in one point, the geometric cancellation criterion
\cite[Theorem~5.4]{Milnor1965HCobordism}.
Absorbing each leaf gives a single ball and carries the non-tree
feet as disjoint disks on its boundary. Reattach those handles with
their actual feet. This gives the lower and upper
diffeomorphisms \(\alpha:P'\to P\) and \(\beta:Q'\to Q\), where
\(P'\cong\HB{\ell}\) and \(Q'\cong\HB{r}\).

For the explicit markings, let \(N\) be the middle cobordism and let
\(p:\partial P\to\partial_-N\) and \(q:\partial Q\to\partial_+N\) be
the given boundary maps. For an incoming framed tube
\(a_j:S^1\times D^2\to\partial_-N\), use the full transported map
\[
 a_j'=(\partial\alpha)^{-1}p^{-1}a_j .
\]
On reversal, the corresponding rule for each dual tube is pullback by
\((\partial\beta)^{-1}q^{-1}\). In this reversal the dual core in the
outgoing tube of the forward handle gives the new attaching circle
and the outgoing \(D^2\) factor gives its normal framing. Let \(X\)
be \(P'\) with the transported \(2\)-handles attached, and let
\(\eta:\partial X\to\partial_+N\) be its induced outgoing marking.
The closed assembly uses
\[
 \varphi=(\partial\beta)^{-1}q^{-1}\eta:\partial X\to\partial Q'.
\]
On the collared pieces the comparison with \(M\) uses
\(\alpha,\operatorname{id}_N,\beta\). A framed slide giving a trace
diffeomorphism \(s:X_{\rm new}\to X\) replaces \(\varphi\) by
\(\varphi\partial s\), and a change of upper coordinates pulls that
map back on the upper side. Thus the framing and not just the core
curve is transported. Tree contraction gives the stated counts,
and \(\chi(M)=2=1-\ell+n_2-r+1\) finishes.
\end{proof}

We use two facts about arbitrary boundary markings.
Hatcher's theorem gives \(\pi_0\operatorname{Diff}(S^3)\cong
\pi_0 O(4)\), so every boundary diffeomorphism of a standard
four-ball extends to the ball by extending its orthogonal endpoint
and its isotopy on a collar \cite[p.~553]{Hatcher1983Smale}.
Laudenbach--Po\'enaru's Theorem A says that
\(\TB{k}\cup_\varphi\HB{k}\cong S^4\) for any boundary
diffeomorphism \(\varphi\). We also use the boundary-extension
statement that a self-diffeomorphism of \(Y_k\) extends over
\(\HB{k}\) \cite[Theorem A and its proof, pp.~337, 341--342]{LaudenbachPoenaru1972Handlebodies}.
An orientation-reversing discrepancy can be absorbed by a
reflection of the rounded product of a three-dimensional handlebody
with an interval. The explicit extension theorem also has an
independent proof \cite[Theorem~2.5]{Moussard2026HandlebodyDiffeomorphisms}.

\begin{proposition}[Zero or one index-two point]\label{prop:one}
If any Morse function on \(M\) has \(n_2\leq1\), then
\(M\cong S^4\).
\end{proposition}

\begin{proof}
For \(n_2=0\), Lemma~\ref{lem:trees} gives two actual balls joined
by their boundary diffeomorphism, which extends by Hatcher's theorem.
For \(n_2=1\), reverse if needed so that \((\ell,r)=(0,1)\).
Then \(X\) is the trace of integer-framed surgery on a knot \(K\) in
\(S^3=\partial D^4\), and the actual boundary is
\(Y_1=S^1\times S^2\), from the upper \(\HB{1}\).
In the exterior of \(K\) its meridian generates \(H_1=\mathbb Z\)
and its preferred longitude vanishes. Filling with integer framing
\(n\) kills \(n\) times the meridian, so \(H_1(Y_1)=\mathbb Z\)
forces \(n=0\). Gabai's Corollary 8.3 computes the knot genus from
the least genus of a nonseparating surface in the zero surgery.
The sphere factor in \(Y_1\) makes the genus zero, so \(K\) is the
unknot \cite[Corollary 8.3 and Remark 8.5]{Gabai1987FoliationsIII}.
Its zero trace is \(\TB{1}\), seen by capping a spanning disk by
the handle core with product normal framing. Theorem A applies to
the given gluing with \(\HB{1}\).
\end{proof}

For comparison, Gabai also gives Po\'enaru's consequence that a
ball with one \(2\)-handle and one \(3\)-handle and with \(H_2=0\)
is a standard ball \cite[Corollary 8.6]{Gabai1987FoliationsIII}.
Scharlemann gives a PL version with one \(2\)-handle and no
\(3\)-handles and uses an end-tree reduction in its proof
\cite[Corollary~2.2 and pp.~422--423]{Scharlemann2008GeneralizedPropertyR}.
The proof of Proposition~\ref{prop:one} uses smooth marked gluings.

\section{Indexed normal area and the middle energy range}\label{sec:index}

Let \(F:M\looparrowright\mathbb R^5\) be a smooth immersion,
\(g=F^*g_{\rm Eucl}\), and \(\nu:M\to S^4\) the oriented unit normal
with \(d\nu=-dF\circ A\). For \(U\subset M\) write
\begin{gather*}
 E_\pm(g;U)=\int_U |W^\pm_g|_{\mathrm{HS}}^2\,dV_g,\qquad
 \Omega_2=\{A\text{ has signature }(2,2)\},\\
 I_2=\int_{\Omega_2}\det A\,dV_g,\qquad
 \omega=|S^4|=\frac{8\pi^2}{3}.
\end{gather*}
Here \(W^\pm_g\) are real symmetric trace-free curvature
endomorphisms on \(\Lambda^\pm\), and unqualified \(E_\pm(g)\)
integrate over \(M\). They agree by signature. All counts and
integrals are on the immersed domain, so distinct preimages at a
self-intersection are distinct points.

\begin{proposition}[Two height averages]\label{prop:averages}
For almost every \(u\in S^4\), \(h_u=\langle F,u\rangle\) is Morse.
If \(n_i(u)\) are its original counts and
\(J=\int_M(-\det A)_+\,dV_g\), then
\begin{align*}
 \frac1\omega\int_{S^4}\sum_{i=0}^4 n_i(u)\,d\sigma
       &=2+\frac{4J}{\omega},\\
 \frac1\omega\int_{S^4}n_2(u)\,d\sigma
       &=\frac{2I_2}{\omega}.
\end{align*}
In particular \(J<\omega\) recognizes the standard smooth sphere.
\end{proposition}

\begin{proof}
By Sard, for almost every direction both \(u\) and \(-u\) are
regular values of \(\nu\). The height is critical precisely on
\(\nu^{-1}(\{u,-u\})\), where its Hessian is
\(\langle u,\nu\rangle g(A\cdot,\cdot)\). This gives nondegeneracy
of all critical points, including for self-intersecting immersions.
Distinctness of their critical values is a separate condition.
The Hessian determinant has the
sign of \(\det A\), and in four dimensions sign reversal leaves the
middle index unchanged. The signed Gauss Jacobian is \(\det A\).
The Morse Euler sum therefore gives \(2\deg\nu=\chi(M)=2\).
The signed area identity yields \(\int_M\det A=\omega\), and
the absolute area identity gives the total count in the display.
Applied on \(\Omega_2\) with both signs of the normal, the same
identity gives the middle count. This is the hypersurface
specialization of total-curvature averaging
\cite{ChernLashof1957Curvature}.

When \(J<\omega\) the average total is less than six. The count
is even by the Morse Euler identity, so some Morse height has at
most four critical points. To list the possible counts, let \(r_j\)
be the rank of the cellular boundary in degree \(j\), with
\(r_0=r_5=0\). Then \(n_j=b_j+r_j+r_{j+1}\). Here the endpoint
Betti numbers are one and the others zero. Two points give
\((1,0,0,0,1)\). Four points give exactly one extra adjacent pair,
so \(n_2\leq1\). Proposition~\ref{prop:one} applies.
\end{proof}

\begin{lemma}[Middle-signature Weyl cost]\label{lem:weyl}
On \(\Omega_2\) one has
\[
 |W^\pm_g|_{\mathrm{HS}}^2\geq \tfrac83\det A.
\]
If a domain contains a nonempty open subset of \(\Omega_2\), the
integrated comparison on that domain is strict.
\end{lemma}

\begin{proof}
Write the principal curvatures as \((-a,-b,c,d)\) with
\(a,b,c,d>0\), and put \(P=abcd\).
The Euclidean Gauss equation diagonalizes the curvature operator
on principal two-forms with entries \(k_i k_j\). Pairing opposite
indices in a chiral basis gives the three sums
\[
 s_1=ab+cd,\qquad s_2=-ac-bd,\qquad s_3=-ad-bc.
\]
The Weyl eigenvalues in each chiral block are \((s_j-\bar s)/2\):
the sectional-pair averages are \(s_j/2\) and the scalar part is
\(\bar s/2\). Hence the squared norm is
\(\frac1{12}\sum_{i<j}(s_i-s_j)^2\). The arithmetic--geometric mean
inequality gives \(s_1\ge2\sqrt P\) and \(s_2,s_3\le-2\sqrt P\),
which proves the comparison. The use of middle-signature
determinants in hypersurface Weyl estimates has the antecedent
\cite[\S\S2--3]{OntiVlachos2017Hypersurfaces}.

Pointwise equality requires \(ab=cd\), \(ac=bd\), and \(ad=bc\),
so \(a=b=c=d\). If this holds on an open set, \(A\) has smooth
rank-two eigenbundles with eigenvalues \(k,-k\), \(k>0\).
For orthogonal unit fields \(X,Y\) in one eigenbundle with
eigenvalue \(\lambda\), Euclidean Codazzi gives
\[
 X(\lambda)=\langle(\nabla_X A)Y,Y\rangle
            =\langle(\nabla_Y A)X,Y\rangle=0 .
\]
Thus \(dk=0\) along both bundles. With \(X\) in the positive
bundle and \(Y\) in the negative bundle, Codazzi now yields
\(-2k(\nabla_XY)_+=2k(\nabla_YX)_-\), so both sides vanish.
Metric compatibility makes each bundle parallel in all directions.
Its mixed sectional curvature is then zero, while Gauss gives
\(-k^2\). The nonnegative continuous defect therefore has positive
integral on any domain containing an open middle-signature subset.
\end{proof}

\begin{theorem}[Recognition and direction selection]\label{thm:band}
The condition \(E_+(g)\le64\pi^2/9\) implies \(M\cong S^4\).
On a nonstandard \(M\), if \(E_+(g)<8\pi^2\), more than three
quarters of normalized spherical directions are Morse heights
with exactly two index-two points. Four mutually orthogonal
directions can be chosen with that count.
\end{theorem}

\begin{proof}
On a nonstandard \(M\), Proposition~\ref{prop:one} gives
\(n_2(u)\ge2\). Thus \(I_2\ge\omega\) by
Proposition~\ref{prop:averages}. Lemma~\ref{lem:weyl} is strict
after integration since \(\Omega_2\) has positive mass, so
\(E_+(g)>\frac83 I_2\ge64\pi^2/9\).
If \(E_+(g)<8\pi^2\), the middle count average is less than \(9/4\).
The proportion of directions with count at least three is bounded
by the average of \(n_2-2\), and is therefore less than \(1/4\).
The first four entries of a random orthonormal frame each have
uniform marginal direction. The union of their four failure events
has probability less than one.
\end{proof}

The two Weyl estimates use the given induced Riemannian metric
without a Yamabe sign hypothesis. Under a smooth conformal change
the chiral Hilbert--Schmidt norm squared and the volume density
scale reciprocally, so the energy condition also tests that
induced conformal class. The count estimate fixes an original
function \(h_u\). It does not specify the distribution of the
other indices, the tree choices, or the framed attaching curves.

\section{Recognized marked surgery presentations}\label{sec:surgery}

Fix a Morse function with exactly two index-two points and apply
Lemma~\ref{lem:trees}. Its possible reduced odd-handle distributions
are \((\ell,r)=(0,2),(1,1),(2,0)\). The following tests are purely
smooth topology and also apply to a marked presentation obtained
without a hypersurface immersion.

\subsection{A one-sided two-component link}

\begin{proposition}[One-sided recognition]\label{prop:onesided}
When \((\ell,r)=(0,2)\), let \(L\) be the actual two-component
framed attaching link in \(\partial D^4\). The same statement
applies after reversal when \((\ell,r)=(2,0)\). If \(L\), possibly
after framed slides, contains an unknotted component or has tunnel
number one, then \(M\cong S^4\).
\end{proposition}

\begin{proof}
The outgoing boundary of its trace is \(Y_2\), by its actual
identification with \(\partial\HB{2}\). In a link exterior the two
meridians generate \(H_1\) freely. In this basis the preferred
longitude of one component represents its linking number with
the other times the other meridian. Thus the two filling slopes
give the integral matrix with framings on the diagonal and
linking number off diagonal. Its cokernel is \(H_1(Y_2)=\mathbb Z^2\),
so it is the zero matrix.

A framed slide replaces an attaching circle by a band sum with
a push-off specified by the other handle's framing. The band
is embedded in the marked three-dimensional level with interior
disjoint from the link. For oriented components in \(S^3\), the
new framing is the sum of the framings plus twice the signed
linking term, as follows by taking the corresponding framed
push-off of the band sum. The slide is an isomorphism of the
smooth trace, with its outgoing marking carried as in
Lemma~\ref{lem:trees}.

For the unknot case, Gompf--Scharlemann--Thompson's
Proposition 3.2 supplies slides to the zero-framed unlink.
It assumes precisely the two-component surgery to \(Y_2\) and
an unknotted component. The proof isotopes the other component
in the zero surgery on the unknot and realizes passages through
the surgery core as slides. For tunnel number one, the tunnel
complement is a genus-two handlebody. Its free group of rank two
surjects onto the link group, which surjects onto \(\pi_1(Y_2)\),
also free of rank two. Hopficity gives a free rank-two link group;
their Proposition 3.1 uses the classical unlink criterion for
this group to identify the unlink
\cite[\S2 and Propositions 2.2, 3.1--3.2]{GompfScharlemannThompson2010}.
Both routes give the trace \(\TB{2}\): split framed-unknot
attachments give the boundary connected sum of their zero traces.
Theorem A applies to the transported gluing with \(\HB{2}\).
\end{proof}

In particular the proof uses the actual link and allows known
framed moves on it. The zero matrix holds for every such link.
Slide equivalence to the unlink for an arbitrary two-component
link with this outgoing boundary is the general Property \(2R\)
question, which is not a premise of Proposition~\ref{prop:onesided}.

\subsection{The two intermediate boundaries in a mixed presentation}

Assume \((\ell,r)=(1,1)\). Denote the marked lower belt sphere
by \(S\subset\partial\HB{1}=S^1\times S^2\) and the two lower
attaching tubes by \(K_1,K_2\), using the same notation for their
cores. In a Kirby diagram let \(U\) be the dotted unknot of the
lower \(1\)-handle. Replacing the dot by a zero coefficient
describes the boundary surgery of the trace by ordinary framed
link surgery in \(S^3\). As a four-dimensional operation this
substitutes \(D^2\times S^2\) for an interior \(S^1\times D^3\);
it changes the trace while retaining the outer boundary. These
statements follow from the product neighborhoods and the
identification of the dotted-unknot boundary with zero surgery
on that unknot. For the dotted-circle convention, handle slides,
and geometric cancellation see
\cite[\S1, pp.~898--899]{Stipsicz2003SpinCobordism}.

The meridian of \(U\) generates \(H_1(\partial\HB{1})\).
Consequently the lower \(2\)-to-\(1\) handle incidence row is
\((\lk(U,K_1),\lk(U,K_2))\). Handles of indices three and four
do not change \(H_1\), so the cokernel of this row is zero.
Orientations, permutation, and actual framed slides of the two
\(2\)-handles implement the integer Euclidean algorithm. Choose
this marked diagram with row \((1,0)\). For the full ordinary
surgery the matrix then has the form
\[
 \begin{pmatrix}0&1&0\\1&d&b\\0&b&c\end{pmatrix}.
\]
Its outgoing boundary is \(Y_1\), so it has rank two and
determinant zero, giving \(c=0\). The intermediate boundaries
have two distinct algebraic possibilities:
\begin{enumerate}
\item Attaching \(K_1\) first gives an integral homology sphere,
by the determinant \(-1\) of its dotted-to-zero two-component minor
and three-dimensional Poincar\'e duality.
\item Attaching \(K_2\) first gives first homology \(\mathbb Z^2\),
since its two-component minor is zero.
\end{enumerate}
The change \(K_2\mapsto K_2-bU\) diagonalizes the remaining
off-block entry of the ordinary surgery matrix. This is a basis
calculation in the ordinary framed surgery trace. It supplies no
geometric cancellation in the original dotted diagram. Keeping just
\(U_0\cup K_2\) yields the \emph{intermediate} boundary in the
second order, while the full ordinary surgery has three components
and boundary \(Y_1\).

\begin{proposition}[A geometric belt intersection]\label{prop:belt}
If after framed slides or isotopy one of the two actual attaching
circles meets \(S\) transversely once, \(M\cong S^4\). The same
criterion holds at the upper end in the dual presentation.
\end{proposition}

\begin{proof}
Attach the indicated handle first. Its attaching circle meets the
belt sphere of the \(1\)-handle in exactly one point of the common
level. Milnor's first cancellation theorem removes the two critical
points by changing the function away from its end levels. Product
flow for this function, fixed in the incoming collar coordinates,
induces a map to the outgoing level. Pull back the other attaching
tube by that map. This product identification sets the outgoing
marking through its flow. In the dotted diagram the same procedure
slides each passage of the other tube
over the cancelling handle to clear the \(1\)-handle disk before
deleting the pair. The bands use disjoint parallel copies along
the cancelling tube with the given framing. In either description
the cancellation changes the marking by the trace map. The
remaining counts are \((1,0,1,1,1)\), or their reverse, and
Proposition~\ref{prop:one} applies. Signed intersection one by
itself allows multiple geometric passages.
\end{proof}

\begin{theorem}[Mixed recognition at a supplied intermediate cut]\label{thm:mixed}
After a choice of framed slides and order of the two middle handles,
let \(Y\) be the actual intermediate boundary in a mixed marked
presentation. If \(Y\cong S^3\) or \(Y\cong Y_2\), then \(M\cong S^4\).
\end{theorem}

\begin{proof}
The two pieces \(X_-,X_+\) at \(Y\) each consist of \(\HB{1}\)
and a single \(2\)-handle, read from their respective ends. The
upper attaching tube is the dual of the second forward handle,
transported in the actual outgoing boundary coordinates.

\emph{The \(S^3\) cut.}
Cap either piece by a marked standard \(D^4\) using a chosen
boundary diffeomorphism. The dual handle attaches to this cap ball
along a framed knot in its \(S^3\), and surgery gives
\(\partial\HB{1}=Y_1\). The surgery homology forces framing zero,
and Property R forces the unknot as in Proposition~\ref{prop:one}.
Thus the capped piece is \(\TB{1}\cup\HB{1}\) with its
transported boundary gluing, and is standard by Theorem A.

Its cap is an actual smoothly parametrized disk in the resulting
\(S^4\). The complement of such a disk is standard: first shrink
the disk by the isotopy of parametrized embeddings \(x\mapsto j(tx)\)
to a small disk about \(j(0)\). In a coordinate chart the rescaling
of \(j(tx)-j(0)\) by \(t^{-1}\) extends smoothly at \(t=0\) to
the derivative of \(j\). This isotopes the small disk to its
local linear ellipsoid and then a coordinate round disk; ambient
isotopy extension transfers the complement. This is the
parametrized disk application of Palais's extension theorem
\cite[Theorem B]{Palais1960Extending}. Hence both \(X_-\) and
\(X_+\) are standard disks and their given gluing is standard by
the boundary extension for \(S^3\).

\emph{The \(Y_2\) cut.}
For either side draw \(\HB{1}\) with dotted unknot \(U\) and
draw the marked framed attaching circle \(K\). The ordinary
two-component surgery \(U_0\cup K\) gives \(Y_2\), so apply the
proof of Gompf--Scharlemann--Thompson's Proposition 3.2. In zero
surgery on \(U\), that proof isotopes \(K\) in \(S^1\times S^2\)
into a three-ball, where Property R identifies it as the unknot.
Here use that isotopy directly in the actual \(\partial\HB{1}\).
It carries the framing and extends through a collar. The framing
in the ball is zero by \(H_1(Y_2)=\mathbb Z^2\): surgery in the
ball must add a free factor to \(H_1(S^1\times S^2)\).
This proves
\[
 X_\pm\cong (S^1\times D^3)\natural(S^2\times D^2),
\]
with the resulting coordinates on the actual gluing.

In the \(S^2\times D^2\) summand of \(X_-\), replace the interior
product neighborhood of \(S^2\times\{0\}\) by \(D^3\times S^1\).
The remaining annulus is a product collar on \(S^2\times S^1\);
thus the replaced \(X_-\) is \(\HB{2}\) with the same outer
marking. The discrepancy between the given gluing and the product
gluing matching boundary connected summands extends over
\(\HB{2}\). The closed surgered manifold is therefore
\((S^1\times S^3)\#S^4\cong S^1\times S^3\). Here the first
summand is the double of \(S^1\times D^3\), and the second is
the union \((S^1\times D^3)\cup(S^2\times D^2)\) in the
product boundary of \(D^2\times D^3\).

The inverse surgery on the actual core circle \(\gamma\) of the
inserted product recovers \(M\). Removing its open tubular
neighborhood preserves fundamental group by general position for
paths and disks relative to a circle in dimension four. Its
boundary has one circle generator, mapping to the class of
\(\gamma\); filling with \(D^2\times S^2\) kills this class.
Thus \(1=\pi_1(M)=\mathbb Z/\langle[\gamma]\rangle\), and
\([\gamma]=\pm1\). Choosing its orientation gives a free
homotopy to the standard factor circle. Relative parametric
transversality makes the homotopy an isotopy: the zero-tangent
condition has codimension four in the two-dimensional parameter
space, and the equal-time double-point condition has codimension
four in the three-dimensional pair-and-time space. By isotopy
extension and tubular-neighborhood uniqueness the complement of
the framed tube is identified with \(S^1\times D^3\). The original
framing induces a boundary diffeomorphism to \(D^2\times S^2\),
and Theorem A at \(k=1\) handles this actual identification.
\end{proof}

In the normalized diagram an integral homology sphere cut is
\(S^3\) if it is simply connected, by the three-dimensional
Poincar\'e theorem. Simple connectivity of \(M\) alone supplies
no such conclusion from van Kampen: at an integral homology sphere
cut each side is contractible already. In fact its single
\(2\)-handle kills the generator of its single \(1\)-handle group,
since for a dotted-to-zero surgery minor with linking \(n\)
one has determinant \(-n^2=\pm1\). Each side then has trivial
fundamental group and no reduced homology, and so is contractible.
Likewise first homology \(\mathbb Z^2\) is only a necessary
condition for the \(Y_2\) cut.

\section{An exterior height test}\label{sec:exterior}

\begin{proposition}\label{prop:exterior}
Cut \(M=P\cup Q\) along a smooth common boundary. Give \(P\) a
smooth immersion into \(\mathbb R^5\) up to the boundary.
Suppose for almost every \(u\in S^4\) its height agrees on \(P\)
with a smooth Morse function on the given \(M\) having at most
\(k(u)\) index-two points in \(\operatorname{int}Q\), where \(k\)
is measurable. Put \(D=\int_{S^4}(2-k(u))_+\,d\sigma\) in ordinary
spherical area. If \(D>0\) and \(E_+(g;P)\le4D/3\), then
\(M\cong S^4\).
\end{proposition}

\begin{proof}
The normal image of the smooth three-dimensional boundary and its
antipode are null. Off those images the scalar extension, whose
first derivative agrees with the known height at the seam, has
no critical point there. Its exterior count follows from
Propositions~\ref{prop:one} and \ref{prop:averages}: on a nonstandard
source there are at least \((2-k(u))_+\) exterior middle points,
so \(2I_2(P)\ge D\). Because \(D>0\), an open exterior set has
middle signature. Lemma~\ref{lem:weyl} gives the strict inequality
\(E_+(g;P)>4D/3\). Only the known exterior is used for the curvature
calculation. The scalar extensions can vary with direction.
\end{proof}

For example, suppose finitely many disjoint patches agree with the
exterior on their marked boundary collars through the quadratic
graphs
\[
 (x,y)\in\Omega_j\longmapsto
 \bigl(R_j(x,y),\,R_jt_j(|y|^2-|x|^2)/2\bigr),
 \qquad R_j,t_j>0 .
\]
Ambient rigid motions can be included in the marks. The filling
heights then have one middle point in patch \(j\) precisely on
the interior normal image \(N_j\) or its antipode, up to null
sets. Indeed the graph Hessian has signature \((2,2)\) and the
upward unit Gauss map is one-to-one. Its gnomonic coordinate is
\(p=-t_j\operatorname{diag}(-1,-1,1,1)(x,y)\). The round metric
of the chart \(p\mapsto(p,1)/\sqrt{1+|p|^2}\) has one
eigenvalue \((1+|p|^2)^{-2}\) in the radial direction and three
\((1+|p|^2)^{-1}\) in transverse directions. Consequently
\[
 |N_j|=t_j^4\int_{\Omega_j}(1+t_j^2|z|^2)^{-5/2}\,dz,\qquad
 k=\sum_j{\bf 1}_{N_j\cup(-N_j)}.
\]
In particular
\[
 \frac D2=\omega-\sum_j |N_j|+
                 \frac12\int_{S^4}(k-2)_+\,d\sigma .
\]
If the domains are copies of a domain star-shaped about zero in
common vertical and horizontal coordinates, their normal images
are nested via the gnomonic coordinates. Then
\(D/2=\omega-|N_{(1)}|-|N_{(2)}|\) uses the two largest areas,
with absent terms zero. This proposition controls the exterior
energy under supplied scalar fillings of the prescribed germs.
It does not provide a bound on the topology of a generic
two-middle-handle diagram.

\section{A calibration on a single embedded sphere}\label{sec:example}

The following example supplies the actual disk and tube data in
a common height level. It is built on the standard sphere.
In a regular height box use \(w=(x_1,x_2,y)\) and \(v\), and let
\(\chi(w,v)\) be a smooth cutoff equal to one on the \(d\)-ball
and zero off the \(2d\)-ball. For small fixed \(a>0\), put
\[
 \begin{aligned}
 Q(w)&=-|x|^2+y^2,\\
 P_{\tau,a}(w,v)&=Q(w)+v^3+3a^2v-6\tau a^2v\chi(w,v),\\
 \psi_a(w,v)&=(w,P_{0,a}(w,v)),\qquad
 U_{\tau,a}=P_{\tau,a}\circ\psi_a^{-1},\qquad 0\le\tau\le1 .
 \end{aligned}
\]
The last coordinate of \(\psi_a\) has positive derivative
\(3(v^2+a^2)\) and is onto at each \(w\). Thus \(U_{0,a}\) is the
affine height coordinate \(T\), and \(U_{\tau,a}=T\) off a
compact set. On the transition annulus \(D(Q+v^3)\) has a
positive minimum in norm and the cutoff perturbation of the
differential is \(O_d(a^2)\). Outside the cutoff the \(v\)
derivative remains positive. In the core at \(\tau=1\) the
polynomial is \(Q+C_a(v)\), \(C_a(v)=v^3-3a^2v\).
Its only critical points are \((0,a)\) of index two and value
\(-2a^3\) and \((0,-a)\) of index three and value \(2a^3\).
Their Hessians are read directly from the quadratic part and
\(C_a''(v)=6v\). The parameter passes through the usual cubic
birth at \(\tau=1/2\)
\cite[\S3.5, equation (3.5.1)]{BorodzikPowellTeichner2025LinkConcordance}.
The disjoint insertions can be carried out at separate parameter
times to give a path with one birth at a time.

Choose two fixed disjoint angular boxes
\(\Theta_j(w)\subset S^3\) on the upper round hemisphere, and take
parameters \(a_2=1.1a_1\) within the above margins. Choose \(T_0\)
so that both insertions and their modified values are supported
in the interior of \((-T_0,T_0)\), and choose \(R_0^2>4T_0+1\).
For sufficiently small \(r>0\), put \(h_*=\sqrt{1-R_0^2r^2}\)
and use the boxes
\[
 z=r\sqrt{R_0^2-2h_*T-r^2T^2}\,\Theta_j(w).
\]
The round graph height at \(z\) is \(h_*+r^2T\). Replace it by
\(h_*+r^2U_{\tau,a_j}(w,T)\), retaining the horizontal coordinates.
The boxes are disjoint and match the round graph on neighborhoods
of their faces. Their heights remain positive, so the horizontal
projection gives an embedding on the upper graph with no collision
with the lower hemisphere. In normalized coordinates \(z/r\)
the latitude and angular inverses have bounded first and second
derivatives on the fixed supporting boxes. Thus the changed graphs
have slopes \(O(r)\) and uniformly bounded Hessians, with the
two fixed birth parameters held constant. Gauss gives
\(|W^\pm|_{\mathrm{HS}}^2\le C|A|^4\) and the changed volume
is \(O(r^4)\). Therefore \(E_\pm=O(r^4)\), uniformly in \(\tau\).
At \(\tau=1\) the actual vertical height has counts
\((1,0,2,2,1)\), since its index-two values are below all its
index-three values and the poles remain the two extrema.

Normalize heights by subtracting \(h_*\) and dividing by \(r^2\).
Set \(d_0=3a_2^3\). The levels \(-d_0,0,d_0\) are regular.
On the relevant local pieces choose the negative coordinate
gradient of \(Q+C_a\) as the descending field, and extend to a
global descending gradient-like field by a partition of unity.
Both fields strictly decrease height on the overlaps away from
critical points. In the box with parameter \(a\), the actual
index-two attaching circle at \(-d_0\) is
\[
 K=\{v=a,\ y=0,\ |x|^2=d_0-2a^3\}.
\]
Near \(v=a\) use \(V=(v-a)\sqrt{v+2a}\) so that
\(C_a(v)=-2a^3+V^2\). This gives the framed tube with
\(x=\sqrt{d_0-2a^3+y^2+V^2}\,\theta\) in the \((y,V)\)
normal coordinates. The same lower level contains the spanning
disk
\[
 \{y=0,\ v_0\le v\le a,\ |x|^2=C_a(v)+d_0\},
\]
where \(v_0\in(-3a,-2a)\) is the simple left root. Its existence
uses \(2<d_0/a^3<4\), since \(C_a(-3a)=-18a^3\) and
\(C_a(-2a)=-2a^3\). The circles collapse smoothly at \(v_0\),
as one can solve for \(v\) in terms of \(|x|^2\) there. The
radial derivative vanishes at the boundary \(v=a\). The vector
\(\partial_y\) is tangent to the level along the disk and is a
disk-normal framing vector, so the handle push-off has zero
linking relative to the disk. Both disks and their tubes lie
in the uncut polynomial regions for the fixed sufficiently
small \(a_j\). The entire lower level is an \(S^3\) by flow from
an unchanged lower latitude. The disks lie in separate angular
boxes, and their three-dimensional regular neighborhoods are
disjoint balls. Thus the actual link in the common incoming
level is the split unlink with both framings zero.

At the actual middle level \(0\) the belt circle of the
index-two handle and the attaching sphere of its index-three
partner in that box are respectively
\[
 \begin{split}
 B&=\{x=0,\ 0\le v\le\sqrt3 a,\ y^2=-C_a(v)\},\\
 T&=\{y=0,\ -\sqrt3 a\le v\le0,\ |x|^2=C_a(v)\}.
 \end{split}
\]
They are smooth at each collapse by simplicity of the relevant
roots. The descending field carries \(B\) toward \((0,a)\) and
the ascending field carries \(T\) toward \((0,-a)\), with both
trajectories confined to the polynomial core. At \(v=0,x=y=0\)
the derivative \(C_a'(0)\ne0\) gives level coordinates \((x,y)\).
The two tangent spaces there are the \(y\)-line and the
\(x\)-plane, so they meet transversely once. The normal line
of the \(3\)-handle tubing is \(\partial_y\) along \(T\).
The two pairs in different boxes are disjoint. The outer levels
of the two boxes give the original latitude product outside
the insertions, and regular-height flow transports its markings.
This verifies both middle handles on one actual height without
needing to reorder index blocks. The construction provides a
standard-type check of the one-sided test.

\section*{Provenance and assistance}
Prepared with OpenAI Codex assistance for Edward B. Baker III on
25 November 2026. The manuscript consolidates the master
immersion section and packets E26, E33, E34, E36, E42, E43, E49,
and E50 at repository base \texttt{483ad3cef5a0af4328c5d35cabbc1f5224d26b1e}.
The manuscript directory records the input revisions and source formats.
The requested execution routing was GPT-6.2 Sol/ultra. Actual
serving model and reasoning effort were unavailable. The manuscript
assembles conditional implications and a checked standard-type
example.

\section{Outstanding challenges}\label{sec:open}

Universal existence of a smooth immersion is a qualitative input:
stable tangent triviality of homotopy spheres, in the rank-five
range on a four-complex, gives a tangent monomorphism into a
trivial rank-five bundle, and the positive-codimension immersion
theorem realizes it \cite[Theorem 3.1]{KervaireMilnor1963Groups}
\cite[\S6]{Hirsch1959Immersions}. Constructing an induced class
below \(8\pi^2\) on every smooth homotopy four-sphere is a
quantitative problem. Shrinking an ambient model by dilation leaves
its four-dimensional curvature energy unchanged. The small energy
in Section~\ref{sec:example} instead uses a bounded Hessian and
a round retained exterior on the standard type.

For a supplied subcritical immersion on a putative nonstandard
type, Theorem~\ref{thm:band} supplies many directions of the right
count, so one can ask for a recognized diagram for even one of
them. A useful selection theorem would have to relate the immersed
geometry of the selected height to its framed tubes on the smooth
domain, using the actual comparison if the critical levels are
reordered. In the mixed case it would suffice to give a belt
intersection as in Proposition~\ref{prop:belt} or to identify the
actual intermediate three-manifold in Theorem~\ref{thm:mixed}.
In the one-sided case it would suffice to find a geometric link
restriction such as that of Proposition~\ref{prop:onesided}, or
otherwise identify the marked trace. The averaged count yields
none of those restrictions. Open coverage of these diagrams and
quantitative subcritical entry are separate dependencies in the
recognition program.

\begingroup
\raggedright
\bibliographystyle{unsrturl}
\bibliography{references}
\endgroup
\end{document}
